\begin{helper}
Let $\mu$ be the push-forward law of the activation--priority output set in
\noderef{RandomIndependentSetSampling}, and let $\nu$ be the underlying law on
pairs $(A,\pi)$.  For every fixed finite set of vertices $Q$,
\[
\operatorname{ofReal}\!\left(\sum_S
  {\bf 1}_{Q\subseteq S}\,\mu(S)\right)
=
\nu\!\left(\{(A,\pi):Q\subseteq I(A,\pi)\}\right),
\]
where
\[
I(A,\pi)=\{z\in A:\pi(w)<\pi(z)\text{ for every }w\in A\cap N_G(z)\}.
\]
\end{helper}
\begin{proof}
The activation coordinate has the discrete sigma algebra, so every test
$z\in A$ is measurable. Each priority evaluation is measurable for the
Borel product sigma algebra, hence so is each strict comparison
$\pi(w)<\pi(z)$. Finite intersections and complements show that
$\{z\in I(A,\pi)\}$ is measurable. For a fixed $S$, its output fiber is the
intersection, over all vertices $z$, of this event when $z\in S$ and its
complement otherwise. Thus every output fiber is measurable; the
survival event below is a finite union of these measurable fibers.

Let $I(A,\pi)$ denote the output set appearing in
\noderef{RandomIndependentSetSampling}.  Put
\[
\mathcal T=\{S:S\text{ is a finite vertex set and }Q\subseteq S\}.
\]
Since the weights $\mu(S)$ are nonnegative, passing the finite sum
\[
\sum_S {\bf 1}_{Q\subseteq S}\mu(S)
\]
through $\operatorname{ofReal}$ gives
\[
\sum_{S\in\mathcal T}\operatorname{ofReal}(\mu(S)).
\]
For each finite set $S$, the push-forward clause of
\noderef{RandomIndependentSetSampling} identifies the summand
\(\operatorname{ofReal}(\mu(S))\) with
\[
\nu\{(A,\pi):S=I(A,\pi)\}.
\]
The events $\{S=I(A,\pi)\}$ are pairwise disjoint as $S$ ranges over distinct
finite sets.  Their union over all $S\in\mathcal T$ is exactly the event
\(\{(A,\pi):Q\subseteq I(A,\pi)\}\): if $Q\subseteq I(A,\pi)$, choose
$S=I(A,\pi)$, and conversely if $S=I(A,\pi)$ with $Q\subseteq S$, then
$Q\subseteq I(A,\pi)$.  Finite additivity of $\nu$ over this disjoint union
therefore gives the displayed identity.
\end{proof}
